\chapter{复现说明}

\section{文档真源}

报告的 LaTeX 真源、编译入口和图件目录分别为：
\begin{itemize}
  \item \path{_paper/technical_report/src/}
  \item \path{_paper/technical_report/build_report.sh}
  \item \path{_paper/technical_report/figures/}
\end{itemize}

各赛道的定量结论来自已归档的 JSON 指标文件。模型、图件与来源之间的对应关系记录在
\path{_paper/technical_report/report_manifest.yml}。后续扩展应先更新赛道实验，
再更新清单与正文，避免手工改写数值造成不一致。

\section{图件生成}

模型架构图采用 Nano Banana 2 生成。初版与审阅修订提示词分别保存在：
\begin{itemize}
  \item \path{_paper/technical_report/prompts/nano_banana_v1/}
  \item \path{_paper/technical_report/prompts/nano_banana_v2/}
  \item \path{_paper/technical_report/prompts/nano_banana_v3/} 至
        \path{_paper/technical_report/prompts/nano_banana_v5/}
\end{itemize}
架构图用于解释信息流，不作为模型性能证据。数据图由下列流水线统一生成：
\begin{itemize}
  \item \path{_pipelines/05_research_visualization_expansion/}
\end{itemize}
真实输入、证据类型、图片尺寸和哈希记录在：
\begin{itemize}
  \item \path{_outputs/research_visualization_expansion/v1/artifact_manifest.json}
\end{itemize}
正文只引用已经人工检查的成品。

基础模型的强基线、随机初始化、预训练初始化和融合条件记录在
\path{_pipelines/05_research_visualization_expansion/foundation_model_experiment_contract.json}。
该合同约束后续实验的可比性，不把待完成实验写成既有结果。

\section{扩展方式}

后续版本可以增加新模型、消融实验、计算开销和新盲测，但六个赛道的固定六节目录
保持不变。新增结果应同时保存数据划分、参数、随机种子、指标文件和图件；若评价集
已经被用于模型选择，必须在正文中继续标明为已知留出集。
